%
\vspace{-0.1cm}
\section{Asynchronous RL Framework}
\label{sec:method}

%
%
%
%

We now present multi-threaded asynchronous variants of one-step Sarsa, one-step Q-learning, n-step Q-learning, and advantage actor-critic.
The aim in designing these methods was to find RL algorithms that can train deep neural network policies reliably and without large resource requirements.
While the underlying RL methods are quite different, with actor-critic being an on-policy policy search method and Q-learning being an off-policy value-based method, we use two main ideas to make all four algorithms practical given our design goal.

%
First, we use asynchronous actor-learners, similarly to the Gorila framework~\citep{nair2015gorila}, but instead of using separate machines and a parameter server, we use multiple CPU threads on a single machine.
Keeping the learners on a single machine removes the communication costs of sending gradients and parameters and enables us to use Hogwild!~\citep{recht2011hogwild} style updates for training.

Second, we make the observation that multiple actors-learners running in parallel are likely to be exploring different parts of the environment.
Moreover, one can explicitly use different exploration policies in each actor-learner to maximize this diversity.
By running different exploration policies in different threads, the overall changes being made to the parameters by multiple actor-learners applying online updates in parallel are likely to be less correlated in time than a single agent applying online updates.
Hence, we do not use a replay memory and rely on parallel actors employing different exploration policies to perform the stabilizing role undertaken by experience replay in the DQN training algorithm.

In addition to stabilizing learning, using multiple parallel actor-learners has multiple practical benefits.
First, we obtain a reduction in training time that is roughly linear in the number of parallel actor-learners.
Second, since we no longer rely on experience replay for stabilizing learning we are able to use on-policy reinforcement learning methods such as Sarsa and actor-critic to train neural networks in a stable way.
%
We now describe our variants of one-step Q-learning, one-step Sarsa, n-step Q-learning and advantage actor-critic.
%

%

\begin{algorithm}[t]
\caption{Asynchronous one-step Q-learning - pseudocode for each actor-learner thread.}
\begin{algorithmic}
\small
\State \emph{// Assume global shared $\theta$, $\theta^-$, and counter $T=0$.}
%
\State Initialize thread step counter $t\gets 0$
\State Initialize target network weights $\theta^- \gets \theta$
\State Initialize network gradients $d\theta \gets 0$
\State Get initial state $s$
\Repeat
\State Take action $a$ with $\epsilon$-greedy policy based on $Q(s,a;\theta)$
\State Receive new state $s'$ and reward $r$
\State $y =
    \left\{
    \begin{array}{l l}
      r  \quad & \text{for terminal } s'\\
      r + \gamma \max_{a'} Q(s', a'; \theta^-) \quad & \text{for non-terminal } s'
    \end{array} \right.$
\State Accumulate gradients wrt $\theta$: $d\theta \gets d\theta + \frac{\partial\left(y - Q(s,a;\theta)\right)^2}{\partial \theta}$
\State $s=s'$
\State $T \gets T + 1$ and $t \gets t + 1$
\If {$T\mod I_{target} == 0$}
\State Update the target network $\theta^- \gets \theta$ 
\EndIf
\If {$t\mod I_{AsyncUpdate} == 0$ or $s$ is terminal}
\State Perform asynchronous update of $\theta$ using $d\theta$.
\State Clear gradients $d\theta \gets 0$.
\EndIf
\Until $T > T_{max}$
\end{algorithmic}
\label{alg:targetq}
\end{algorithm}

\textbf{Asynchronous one-step Q-learning:}
Pseudocode for our variant of Q-learning, which we call Asynchronous one-step Q-learning, is shown in Algorithm~\ref{alg:targetq}.
Each thread interacts with its own copy of the environment and at each step computes a gradient of the Q-learning loss.
%
We use a shared and slowly changing target network in computing the Q-learning loss, as was proposed in the DQN training method.
%
We also accumulate gradients over multiple timesteps before they are applied, which is similar to using minibatches.
This reduces the chances of multiple actor learners overwriting each other's updates.
%
Accumulating updates over several steps also provides some ability to trade off computational efficiency for data efficiency.

%
Finally, we found that giving each thread a different exploration policy helps improve robustness.
Adding diversity to exploration in this manner also generally improves performance through better exploration.
While there are many possible ways of making the exploration policies differ we experiment with using $\epsilon$-greedy exploration with $\epsilon$ periodically sampled from some distribution by each thread.

%
%
%
%

%

\textbf{Asynchronous one-step Sarsa:}
The asynchronous one-step Sarsa algorithm is the same as asynchronous one-step Q-learning as given in Algorithm~\ref{alg:targetq} except that it uses a different target value for $Q(s,a)$.
The target value used by one-step Sarsa is $r + \gamma Q(s', a'; \theta^-)$ where $a'$ is the action taken in state $s'$~\citep{rummery1994sarsa,sutton:book}.
%
%
%
%
%
%
%
%
%
We again use a target network and updates accumulated over multiple timesteps to stabilize learning.

%

%
%
%
%
%
%
%
%
%
%
%
%
%
%
%
%
%
%
%
%
%
%
%
%
%
%
%
%
%
%
%
%
%
%
%
%
%
%
%
%
%

\textbf{Asynchronous n-step Q-learning:}
Pseudocode for our variant of multi-step Q-learning is shown in Supplementary Algorithm~\ref{alg:msq}.
The algorithm is somewhat unusual because it operates in the forward view by explicitly computing n-step returns, as opposed to the more common backward view used by techniques like eligibility traces~\citep{sutton:book}.
We found that using the forward view is easier when training neural networks with momentum-based methods and backpropagation through time.
In order to compute a single update, the algorithm first selects actions using its exploration policy for up to $t_{max}$ steps or until a terminal state is reached.
This process results in the agent receiving up to $t_{max}$ rewards from the environment since its last update.
The algorithm then computes gradients for n-step Q-learning updates for each of the state-action pairs encountered since the last update.
Each n-step update uses the longest possible n-step return resulting in a one-step update for the last state, a two-step update for the second last state, and so on for a total of up to $t_{max}$ updates.
The accumulated updates are applied in a single gradient step.

%

%
%
%
%
%
%
%
%
%
%
%
%
%
%
%
%
%
%
%
%
%
%
%
%
%
%
%
%
%
%
%
%
%
%
%
%
%

\textbf{Asynchronous advantage actor-critic:}
The algorithm, which we call asynchronous advantage actor-critic (A3C), maintains a policy $\pi(a_t|s_t;\theta)$ and an estimate of the value function $V(s_t;\theta_v)$.
Like our variant of n-step Q-learning, our variant of actor-critic also operates in the forward view and uses the same mix of n-step returns to update both the policy and the value-function.
The policy and the value function are updated after every $t_{max}$ actions or when a terminal state is reached.
The update performed by the algorithm can be seen as $\nabla_{\theta'} \log\pi(a_t|s_t;\theta') A(s_t,a_t;\theta,\theta_v)$ where $A(s_t,a_t;\theta,\theta_v)$ is an estimate of the advantage function given by $\sum_{i=0}^{k-1}\gamma^i r_{t+i} + \gamma^k V(s_{t+k};\theta_v)-V(s_t;\theta_v)$, where $k$ can vary from state to state and is upper-bounded by $t_{max}$.
The pseudocode for the algorithm is presented in Supplementary Algorithm~\ref{alg:reinforce}.


As with the value-based methods we rely on parallel actor-learners and accumulated updates for improving training stability.
Note that while the parameters $\theta$ of the policy and $\theta_v$ of the value function are shown as being separate for generality, we always share some of the parameters in practice.
We typically use a convolutional neural network that has one softmax output for the policy $\pi(a_t|s_t;\theta)$ and one linear output for the value function $V(s_t;\theta_v)$, with all non-output layers shared.

We also found that adding the entropy of the policy $\pi$ to the objective function improved exploration by discouraging premature convergence to suboptimal deterministic policies.
This technique was originally proposed by~\citep{williams1991function}, who found that it was particularly helpful on tasks requiring hierarchical behavior.
The gradient of the full objective function including the entropy regularization term with respect to the policy parameters takes the form
%
$\nabla_{\theta'} \log\pi(a_t|s_t;\theta') (R_t - V(s_t;\theta_v)) + \beta \nabla_{\theta'} H(\pi(s_t;\theta'))$,
%
where $H$ is the entropy.  The hyperparameter $\beta$ controls the strength of the entropy regularization term.


%
\textbf{Optimization:}
We investigated three different optimization algorithms in our asynchronous framework -- SGD with momentum, RMSProp~\citep{tieleman2012lecture} without shared statistics, and RMSProp with shared statistics.
We used the standard non-centered RMSProp update given by
\begin{eqnarray}
\label{eq-rmsprop}
g = \alpha g + (1-\alpha)\Delta\theta^2 \text{ and } \theta \gets \theta - \eta \frac{\Delta\theta}{\sqrt{g+\epsilon}},
\end{eqnarray}
where all operations are performed elementwise.
A comparison on a subset of Atari 2600 games showed that a variant of RMSProp where statistics $g$ are shared across threads is considerably more robust than the other two methods.
Full details of the methods and comparisons are included in Supplementary Section~\ref{sec:opt}.
